\begin{answers}{ch:attention}{注意機構（Attention）}
\item $\vy_2=(0.196,\ 0.804)$ で、実行結果と一致します。
2 番目のトークンは 1〜2 番目だけを見ます。
$\vq_2=(0,2)$ と $\vk_1=(1,0)$、$\vk_2=(0,1)$ の内積は $0$ と $2$ です。
$\sqrt2$ で割ると $0$ と $1.414$、指数をとると $1$ と $4.113$ で、合計は $5.113$ です。
重みは $1/5.113\approx0.196$ と $4.113/5.113\approx0.804$ になります。
よって $\vy_2=0.196\,(1,0)+0.804\,(0,1)=(0.196,\ 0.804)$ です。

\item 分母を $Z=\sum_k e^{z_k}$ とおきます。$\pd{Z}{z_j}=e^{z_j}$ です。
$i=j$ のとき、商の微分法より $\pd{a_i}{z_i}=\dfrac{e^{z_i}Z-e^{z_i}e^{z_i}}{Z^2}=a_i-a_i^2=a_i(1-a_i)$ です。
$i\ne j$ のとき、分子 $e^{z_i}$ は $z_j$ を含まないので、$\pd{a_i}{z_j}=\dfrac{-e^{z_i}e^{z_j}}{Z^2}=-a_ia_j$ です。
二つをまとめると $a_i(\delta_{ij}-a_j)$ になります。
$a=(1,0,0)$ では、対角は $1\cdot0$ と $0\cdot1$、対角以外も $-1\cdot0$ や $0\cdot0$ で、9 個すべてが 0 です。
$a=(\frac13,\frac13,\frac13)$ では、対角が $\frac13\cdot\frac23=\frac29\approx0.222$、対角以外が $-\frac19\approx-0.111$ です。
one-hot では勾配が完全に消え、一様に近いほど勾配が残ることがわかります。

\item 正答率は 96.0\%（$H=4$、\tabref{tab:attn-experiment}）から 73.0\% に下がりました（答えの損失は 0.032 から 0.239）。
\code{T, D, H = 12, 64, 4} の \code{H} を 1 にし、注意のモデルだけを同じ 3,000 ステップ学習して測りました（乱数の種 0）。
ただし、この値は乱数の種で大きく変わります。
\figref{fig:attn-sweep-heads}のとおり、種 1 では $H=1$ の正答率は 14.8\% で、「平均」のモデルの 40.2\% より低くなりました。
$H=4$ では 96.0\% と 95.0\%、$H=8$ では 99.6\% と 98.4\% で、ヘッドが多いほど正答率が高く、種による差も小さくなります。
答えの一の位を書くには、二つの数の一の位を両方取ってくる必要があります。
重み付き平均で両方を取ると、二つの数字のバリューが 1 本のベクトルに混ざって届きます。
\secref{sec:attn-visualize}の基盤モデルでは、二つのヘッドが別々の数字を 1 個ずつ見て、混ざらずに運んでいました。
1 ヘッドではこの分担ができないので、正答率が下がり、学習がうまく進むかどうかも種に左右されたと考えられます。
種 0 のように 1 ヘッドでも 73\% 解けることがあるので、混ざった情報からでも、ある程度は後ろの MLP が答えを取り出せるようです。

\item 「\code{=}」の行では、層 3・ヘッド 1 が十の位の「\code{1}」と「\code{3}」に 0.42 ずつ、層 3・ヘッド 3 が「\code{3}」に 0.91 を置いていました。
層 1・ヘッド 0 は前と同じく「\code{=}」に 0.95 です（\code{visualize.py} の \code{text} だけを書き換えて実行）。
答えの十の位を書く位置で十の位の数字を見る点は、\figref{fig:attn-heat-drill}と同じです。
一方、繰り上がりを決める一の位の「\code{2}」「\code{5}」には、この 2 つのヘッドはほとんど重みを置いていません（0.08 以下）。
実際、「\code{=}」の次の予測は「\code{5}」が 0.559、正解の「\code{4}」が 0.428 で、貪欲に生成すると「57」と誤答します。
繰り上がりを確かめずに当て推量している、と読めます（重みの解釈には限界があることに注意）。
なお「\code{4}」の行では、ヘッド 1 が「\code{5}」に 0.94、ヘッド 3 が「\code{2}」に 0.99 と、一の位を分担して見ていました。

\item 約 $4.0\times10^9$ 回で、$4Td^2$ のほうが大きくなります。
$T=1024$、$d=768$ を代入すると、$4Td^2=4\times1024\times768^2\approx2.42\times10^9$、$2T^2d=2\times1024^2\times768\approx1.61\times10^9$ で、合計は約 $4.03\times10^9$ です。
2 つの項の比は $4Td^2/(2T^2d)=2d/T$ です。
これが 1 より小さくなるのは $T>2d=1536$ のときです。
GPT-2 では、文脈長が 1,536 を超えると $T^2$ の項のほうが大きくなります。
\end{answers}
